% MT008 / explanation — canonical editable content.
先读横轴\(z=yf(x)\)。标签\(y\)取\(-1\)或\(+1\)：分数与标签同号时\(z>0\)，分类正确；异号时\(z<0\)，分类错误。\(z\)还表示朝正确方向走了多远，而不只是对错。


图像是一条最低点在\(z=1\)的平方形曲线，可用\((z-1)^2\)理解。它会把每个训练点的有符号分数往1拉。


\textbf{左侧：错得很远的点受很大惩罚。} 以教学公式\((z-1)^2\)算，\(z=-1\)的损失是4，\(z=-4\)时为25；斜率\(2(z-1)\)的大小也从4增加到10。极端错分样本会强烈拉动模型，若它是噪声或错误标签，就可能影响整体拟合。


\textbf{右侧：已经很有把握的正确点仍被罚。} \(z=1\)损失0，但\(z=4\)损失又成9。模型需要把这个正确分数拉低到1。这与hinge在\(z\ge1\)时不再处罚不同：对分类任务来说，已经正确且留有间隔的点未必需要继续调整。


原解据此评价它不理想；实际准确率还取决于数据与模型。


\textbf{出处：} 2020B期中Quiz Q8，原卷第4页（10分）


\textbf{答案依据：} 2020B期中 Q8 随卷／配套参考答案；条件与解释经复核。
